\documentclass[11pt,a4paper]{article}
\usepackage[margin=25mm]{geometry}
\usepackage[T1]{fontenc}
\usepackage[utf8]{inputenc}
\usepackage{lmodern,amsmath,amssymb,enumitem,xurl,xcolor}
\usepackage[colorlinks=true,linkcolor=black,urlcolor=blue!55!black]{hyperref}
\hypersetup{pdfauthor={Luca Scharr},pdftitle={From Images to 3D}}
\setlength{\parindent}{0pt}
\setlength{\parskip}{0.45em}
\setlist{leftmargin=*,itemsep=0.15em,topsep=0.15em}
\emergencystretch=2em
\begin{document}
\begin{center}
{\Large\bfseries Tutorial 10 - From Images to 3D - Worked Solutions}\\[2mm]
Luca Scharr\\
Computer Vision 2026\\
University of Bern\\[1mm]
Interactive lab: \href{https://sosekie.github.io/Unibe-CVG-CV-Interactive-Labs/tutorial-10/}{Open Tutorial 10 online}
\end{center}
Each question is repeated before its worked answers.
\section*{Part 1. One image and more evidence}
\subsection*{1. One image, many scenes}
\par
\noindent
\begin{minipage}{\linewidth}
\subsubsection*{1.1}
A point $(X,Y,Z)$ projects to $x=fX/Z$ with a calibrated pinhole camera.
\begin{enumerate}[label=(\alph*),nosep]
\item Multiply all three point coordinates by $k>0$.
What happens to x?
\item Does an unchanged image establish that the physical scene is unchanged?
\end{enumerate}
\noindent\textbf{Solution (a).} $x'=f(kX)/(kZ)=fX/Z=x$.
The same argument holds for the vertical coordinate.

\noindent\textbf{Solution (b).} No.
The point is on the same viewing ray but its distance has changed.
Images alone do not reveal absolute scale; independent size or distance information can resolve it.
\end{minipage}
\par\medskip
\par
\noindent
\begin{minipage}{\linewidth}
\subsubsection*{1.2}
An object is twice as large and twice as far from the camera.
The camera does not move.
\begin{enumerate}[label=(\alph*),nosep]
\item What happens to its projected size under perspective projection?
\item Would a learned estimate of its size make the observation geometrically unique?
\end{enumerate}
\noindent\textbf{Solution (a).} The ratio of object size to distance stays fixed, so the projected size is unchanged under the stated scaling.

\noindent\textbf{Solution (b).} No.
A learned size prior can favour a plausible interpretation, but an unusual object can violate it.
The input image still admits multiple 3D explanations.
\end{minipage}
\par\medskip
\noindent\href{https://sosekie.github.io/Unibe-CVG-CV-Interactive-Labs/tutorial-10/#ambiguity}{Open the interactive experiment.}
\subsection*{2. Combine a prior and a second view}
\par
\noindent
\begin{minipage}{\linewidth}
\subsubsection*{2.1}
A calibrated camera has focal length f = 1.
An object projects to image height h = 0.5, so $h=fH/Z$, where H is object height and Z is distance.
\begin{enumerate}[label=(\alph*),nosep]
\item Calculate Z for H = 1.5 m and H = 3 m.
Can this image alone distinguish them?
\item Why could a model trained mostly on 1.5 m objects give a plausible but wrong metric estimate for a 3 m object?
\end{enumerate}
\noindent\textbf{Solution (a).} Rearranging gives $Z=fH/h$.
The two candidate distances are 3 m and 6 m.
Both produce the same image height, so that measurement alone does not distinguish them.

\noindent\textbf{Solution (b).} A learned size preference can favour H = 1.5 m and therefore Z = 3 m, although the actual object may be 3 m tall at Z = 6 m.
The prior favours a likely interpretation; it does not add geometric evidence of the true size.
\end{minipage}
\par\medskip
\par
\noindent
\begin{minipage}{\linewidth}
\subsubsection*{2.2}
A toy model assigns probabilities to candidate object heights and uses the fixed image measurement h = 0.5 to obtain candidate distances.
It then receives a disparity measurement from a second calibrated camera.
\begin{enumerate}[label=(\alph*),nosep]
\item What determines the distance distribution before the second measurement?
\item How should the second measurement change the candidate probabilities, and what if that measurement is noisy?
\end{enumerate}
\noindent\textbf{Solution (a).} Each candidate height H maps to $Z=fH/h$.
The assigned size probabilities therefore determine the probabilities of these image-consistent distances.
A broader size prior yields a broader range of distance hypotheses.

\noindent\textbf{Solution (b).} Multiply each prior candidate probability by the likelihood of the measured disparity under that candidate, then normalize.
A noisy measurement gives a broader likelihood and generally a less concentrated update; a biased or incorrect measurement can favour the wrong distance.
\end{minipage}
\par\medskip
\par
\noindent
\begin{minipage}{\linewidth}
\subsubsection*{2.3}
Two parallel calibrated cameras with focal length f and baseline B observe a point with disparity $d=u_1-u_2$.
\begin{enumerate}[label=(\alph*),nosep]
\item Derive its depth Z in terms of $f,B,\text{ and }d$.
\item Why does a small baseline make depth sensitive to matching error?
\end{enumerate}
\noindent\textbf{Solution (a).} For the camera pair, $u_1=fX/Z,\quad u_2=f(X-B)/Z$, so $d=fB/Z,\quad Z=fB/d$.

\noindent\textbf{Solution (b).} For fixed depth, smaller B produces smaller disparity.
Since $\frac{\mathrm d Z}{\mathrm d d}=-fB/d^2$, a fixed disparity error can cause a large depth error when disparity is small.
\end{minipage}
\par\medskip
\par
\noindent
\begin{minipage}{\linewidth}
\subsubsection*{2.4}
A second camera rotates around the same optical centre as the first.
A point is matched accurately across the images.
\begin{enumerate}[label=(\alph*),nosep]
\item Does this rotation alone triangulate the distance of the point?
\item What scale remains ambiguous if camera translation is estimated only from images?
\end{enumerate}
\noindent\textbf{Solution (a).} No.
Both rays share an origin, so a point can lie at any positive distance along the common geometric ray, despite changing image coordinates.

\noindent\textbf{Solution (b).} Without a known baseline or another metric reference, scaling all points and camera translations together preserves the image projections.
Relative geometry can still be constrained.
\end{minipage}
\par\medskip
\par
\noindent
\begin{minipage}{\linewidth}
\subsection*{Optional extension: Shortcut learning}
A toy model predicts depth as $4+w_{\rm size}(1/h-2)+w_{\rm colour}c$, where h is image height and c = -1 for red or +1 for blue.
The training set contains a red object at 2 m with h = 1 and a blue object at 6 m with h = 1/3.
\begin{enumerate}[label=(\alph*),nosep]
\item Ignoring regularisation, what constraint must the weights satisfy to fit both training examples exactly?
Give two different solutions.
\item Now add a blue object at 2 m and a red object at 6 m.
Which weights fit all four examples exactly, and why does that change generalisation to swapped colours?
\end{enumerate}
\noindent\textbf{Solution (a).} For the two original examples, both input features are -1 for the red object and +1 for the blue object.
Both observations therefore give the same constraint: $w_{\rm size}+w_{\rm colour}=2$.
For example, (2, 0) uses image size alone and (0, 2) uses colour alone; both fit the training set.

\noindent\textbf{Solution (b).} The swapped examples add the independent constraint $w_{\rm size}-w_{\rm colour}=2$.
Together the constraints give $w_{\rm size}=2,\quad w_{\rm colour}=0$.
Colour no longer explains the training labels, so a fit can rely on the size feature instead.
The interactive model uses a small weight penalty, so its fitted weights are close to, rather than exactly, these unregularised values.
\end{minipage}
\par\medskip
\noindent\href{https://sosekie.github.io/Unibe-CVG-CV-Interactive-Labs/tutorial-10/#learning}{Open the interactive experiment.}
\section*{Part 2. Fit a scene}
\subsection*{3. Fit structure and motion}
\par
\noindent
\begin{minipage}{\linewidth}
\subsubsection*{3.1}
In the experiment, cameras 1 and 2 are at X = 0 and X = 1 with focal length f = 1.
Landmark A has synthetic reference coordinates $(X,Y,Z)=(-0.65,0.42,3.2)$.
The initial estimated coordinates are (-0.42, 0.53, 2.4).
Each camera projects a point to $(u,v)=((X-C)/Z,Y/Z)$, where C is the camera position.
\begin{enumerate}[label=(\alph*),nosep]
\item Compute the observed and initial projected coordinates for A in each of the first two images, and compute the mean of the two squared 2D residuals.
\item The fitting control also adjusts the other three points and camera 3.
Why must these parameters be considered together in the full objective?
\end{enumerate}
\noindent\textbf{Solution (a).} Camera 1 observes (-0.203125, 0.131250) and initially projects A to (-0.175000, 0.220833), giving squared residual approximately 0.008816.
Camera 2 observes (-0.515625, 0.131250) and initially projects A to (-0.591667, 0.220833), giving approximately 0.013808.
Their mean is approximately 0.011312.
The website metric averages over all visible landmarks, not only these two observations.

\noindent\textbf{Solution (b).} Each 3D point projects into every image in which it is visible, and camera 3 affects the projections of all four points in its image.
The full sum couples the shared scene and camera parameters across observations.
This fixture fixes camera 1, camera 2, their 1 m baseline, focal length and orientations; a general bundle-adjustment problem can estimate more parameters.
\end{minipage}
\par\medskip
\par
\noindent
\begin{minipage}{\linewidth}
\subsubsection*{3.2}
All four landmarks are visible in all three images, giving 12 two-dimensional observations.
The experiment corrupts landmark C in camera 3 by adding 0.22 to its horizontal observation and subtracting 0.12 from its vertical observation.
\begin{enumerate}[label=(\alph*),nosep]
\item If the estimated geometry initially equals the true geometry, calculate the corrupted observation's squared 2D residual and the mean squared residual across all 12 observations.
\item What can happen when the fitting algorithm tries to reduce this error, and how could one limit the influence of the wrong match?
\end{enumerate}
\noindent\textbf{Solution (a).} The residual at the corrupted observation is (-0.22, 0.12), whose squared norm is 0.22 squared plus 0.12 squared, or 0.0628.
With the other 11 residuals equal to zero, the mean squared 2D residual is 0.0628/12, approximately 0.005233.

\noindent\textbf{Solution (b).} The optimiser may move landmark C and shared camera 3 to fit the incorrect observation, thereby distorting geometry or worsening agreement in other views.
Correspondence verification and robust estimation or a robust loss can reduce an outlier's influence.
A low residual by itself does not verify the recovered geometry.
\end{minipage}
\par\medskip
\subsection*{Optional SfM factorisation}
Why does centering remove the image translations?
Why does any invertible $3\times 3$ matrix leave the factorisation ambiguous?

\noindent\textbf{Solution.} Subtracting the image centroid removes each view's translation and gives $D=MS$, so $\operatorname{rank}(D)\leq 3$.
For any invertible $C$, $(MC)(C^{-1}S)=MS$, so the factors are not unique.
Orthographic camera constraints are needed to reduce this affine ambiguity.
The rank can be smaller in degenerate configurations.
\noindent\href{https://sosekie.github.io/Unibe-CVG-CV-Interactive-Labs/tutorial-10/#sfm}{Open the interactive experiment.}
\section*{Part 3. Representations and evaluation}
\subsection*{4. Choose a 3D representation}
\par
\noindent
\begin{minipage}{\linewidth}
\subsubsection*{4.1}
Choose a representation for measuring visible distances, building a closed object mesh, and rendering a nearby novel view.
\begin{enumerate}[label=(\alph*),nosep]
\item Which representation is suitable for each task?
\item What limitation or validation requirement accompanies each choice?
\end{enumerate}
\noindent\textbf{Solution (a).} A calibrated metric depth or point map can support visible distance measurements; an occupancy field can be converted to a mesh; a Gaussian representation can render a nearby view.

\noindent\textbf{Solution (b).} Depth misses occluded surfaces and depends on calibration and scale.
Hidden occupancy needs evidence or a learned prior and geometric validation.
A convincing Gaussian rendering does not itself guarantee a watertight or accurate collision surface.
\end{minipage}
\par\medskip
\par
\noindent
\begin{minipage}{\linewidth}
\subsubsection*{4.2}
A network predicts a depth z and a camera ray direction r for one pixel.
\begin{enumerate}[label=(\alph*),nosep]
\item How do the direction convention and the definition of depth determine the 3D point?
\item Why can two plausible predicted hidden surfaces be consistent with the same visible image?
\end{enumerate}
\noindent\textbf{Solution (a).} For camera centre c, $X=c+zr$ only if r is parameterized so its optical-axis component is one and z is axial depth.
If r is unit length, the multiplier must instead be distance along the ray.

\noindent\textbf{Solution (b).} Occluded geometry does not directly contribute to the visible image.
Appearance and learned priors can favour a completion, but one view may not determine the hidden surface uniquely.
\end{minipage}
\par\medskip
\noindent\href{https://sosekie.github.io/Unibe-CVG-CV-Interactive-Labs/tutorial-10/#representations}{Open the interactive experiment.}
\subsection*{5. Evaluate depth and geometry}
\par
\noindent
\begin{minipage}{\linewidth}
\subsubsection*{5.1}
Prediction A is exactly twice the reference depth at every pixel.
Prediction B predicts the wall depth for every pixel and therefore misses the narrow foreground pole.
\begin{enumerate}[label=(\alph*),nosep]
\item What happens to A's error if a best-fit global scale is permitted?
\item Why can B have a small whole-image error while missing an important object?
\end{enumerate}
\noindent\textbf{Solution (a).} Multiplying A by 0.5 exactly recovers the reference under this idealized example.
Its metric error before alignment is not zero, so the chosen evaluation protocol matters.

\noindent\textbf{Solution (b).} The pole occupies only 16 of 240 pixels.
Its squared error is 16 square metres per pole pixel, but the 224 wall pixels have zero error, so the whole-image mean is approximately 1.067 square metres.
Inspect the pole-region error as well as the average.
\end{minipage}
\par\medskip
\par
\noindent
\begin{minipage}{\linewidth}
\subsubsection*{5.2}
A Gaussian reconstruction matches its training photographs well and yields a plausible held-out rendering.
\begin{enumerate}[label=(\alph*),nosep]
\item What evidence does the held-out rendering provide?
\item Does that alone validate a collision surface in hidden regions?
\end{enumerate}
\noindent\textbf{Solution (a).} It tests novel-view appearance for that camera and scene, beyond the training rays.
It is stronger evidence for view synthesis than training fit alone, subject to test selection.

\noindent\textbf{Solution (b).} No.
Hidden surfaces and distances can be unconstrained even when images look plausible.
Define a surface extraction procedure and compare against independent geometry or task-relevant measurements.
\end{minipage}
\par\medskip
\noindent\href{https://sosekie.github.io/Unibe-CVG-CV-Interactive-Labs/tutorial-10/#evaluation}{Open the interactive experiment.}
\end{document}
